% B12 / A058 proof candidate. Full FEI remains open.
\section*{Sparse quotient-union overlap theorem}

Identify the uniform cube with $\mathbb F_2^n$.  Let
$L:\mathbb F_2^n\to\mathbb F_2^r$ be surjective, let
$B\subseteq\mathbb F_2^r$ be nonempty, and put
\[
 A=L^{-1}(B),\qquad p=|B|2^{-r},\qquad
 K(x,z)=1-2\mathbf 1_{\{x=-1\}}\mathbf 1_A(z).
\]
Thus $K$ is a genuine Boolean one-bit composition between the
constant branch and $D_A=1-2\mathbf1_A$.  The parity rows defining
$L$ may overlap arbitrarily in their physical input coordinates.

\begin{theorem}[Sparse quotient-union overlap]
If $p\le1/4$ and $\ell=\log_2(1/p)\ge r/2$, then
\[
 \mathbf H(K)<8\mathbf I(K).
\]
Equivalently, it is enough that $|B|\le2^{r/2}$ and $p\le1/4$.
The selector influence is exactly $p$, so the theorem has no
uniform positive selector-energy hypothesis.  If $|B|\ge3$, this
includes non-affine unions of quotient cells; the two-cell case is
not new because any two points of $\mathbb F_2^r$ form an affine line.
\end{theorem}

\begin{proof}
Write
\[
 a_\alpha=\mathbb E_z\left[\mathbf1_A(z)
       (-1)^{\alpha\cdot Lz}\right],\qquad \alpha\in\mathbb F_2^r.
\]
Surjectivity identifies these with the normalized Fourier
coefficients of $\mathbf1_B$; hence
\[
 a_0=p,\qquad \sum_\alpha a_\alpha^2=p.
\]
Since
\[
 K(x,z)=1-\mathbf1_A(z)+x\mathbf1_A(z),
\]
its constant Fourier coefficient is $1-p$.  For every nonzero
$\alpha$, the two frequencies $L^T\alpha$ and
$\{x\}\cup L^T\alpha$ have squared coefficient $a_\alpha^2$,
and the selector singleton has squared coefficient $p^2$.
Thus the total nonconstant spectral mass is
\[
 q=p^2+2\sum_{\alpha\ne0}a_\alpha^2=2p-p^2.
\]
There are at most $2^{r+1}-1$ nonconstant atoms.  Conditional
entropy maximization at fixed total mass therefore gives the exact
support ledger
\[
 \mathbf H(K)
 \le h_2(q)+q\log_2(2^{r+1}-1)
 < h_2(q)+q(r+1).                                      \tag{1}
\]

The selector coordinate contributes exactly $p$.  On data
coordinates, the two Fourier sectors have equal squared masses, so
\[
 \mathbf I(K)=p+2\sum_{\alpha\ne0}|L^T\alpha|a_\alpha^2
             =p+\tfrac12\mathbf I(D_A).                \tag{2}
\]
The cube edge-isoperimetric inequality, applied to the set $A$ of
density $p\le1/2$, gives
\[
 \mathbf I(D_A)\ge2p\log_2(1/p)=2p\ell.
\]
Consequently
\[
 \mathbf I(K)\ge p(1+\ell).                            \tag{3}
\]

Because $p\le1/4$, $q\le1/2$.  For $0<q\le1/2$,
$h_2(q)\le2q\log_2(1/q)$; this follows by comparing the two
summands of $h_2(q)$.  Also $p\le q\le2p$.  Therefore (1) implies
\[
 \mathbf H(K)<4p\ell+2p(r+1).                          \tag{4}
\]
Combining (3)--(4),
\[
 \frac{\mathbf H(K)}{\mathbf I(K)}
 <\frac{4\ell+2(r+1)}{\ell+1}
 =4+\frac{2(r-1)}{\ell+1}<8,
\]
where the last inequality uses $\ell\ge r/2$.  This proves the
claim.  Every quantity in (1)--(4) refers to the actual shared
physical inputs; no disjoint-block entropy rule is used.
\end{proof}

\paragraph{Disjoint leaf tensorization.}
If a Boolean leaf $w$ uses variables disjoint from $(x,z)$ and
$\mathbf H(w)\le C\mathbf I(w)$, then for $F=wK$ the Fourier law is
a product and
\[
 \mathbf H(F)=\mathbf H(w)+\mathbf H(K),\qquad
 \mathbf I(F)=\mathbf I(w)+\mathbf I(K).
\]
Hence $\mathbf H(F)<\max\{C,8\}\mathbf I(F)$.  This final step is
explicitly disjoint tensorization; it is not a shared-variable
product claim.
